\chapter{Innovation Charter and Opportunity Identification}
\label{sec:innovation-charter}

\section{Session Overview and Objectives}

This session introduces the \textbf{innovation charter}, which is a short statement that defines the direction and boundaries of your project. A good charter does not lock the team into one detailed solution. Instead, it acts as a compass that keeps opportunity generation focused while still leaving room for exploration.

By the end of this session, you should be able to:
\begin{enumerate}
    \item Explain the purpose of an innovation charter in product development;
    \item Draft multiple charter statements with appropriate breadth;
    \item Evaluate charters using relevance, scope, feasibility, and physics-related potential; and
    \item Agree on one final charter for the team to use in Session~\ref{ch:opportunity-identification}.
\end{enumerate}

\section{What Is an Innovation Charter?}

An \textbf{innovation charter} is a concise statement that defines the goals and boundary conditions of an innovation effort. It usually indicates the product category or problem area, the intended market or user context, and the time horizon or strategic intent. It should be broad enough to encourage idea generation, but clear enough to prevent the team from drifting into unrelated topics.

For APID class, the charter should remain connected to hydroponic agriculture and leave space for meaningful applied physics work such as sensing, flow control, structural design, or energy management.

\section{Example Charters}

The example below shows the style of statement you should aim for.

\begin{quote}
``Create a physical product in the cat toy category that we can launch to the market within about a year through our existing retail sales channel.''

\hfill --- FroliCat example
\end{quote}

This charter is effective because it identifies a category, sets a rough time frame, and establishes a market context without prescribing a specific design.

A hydroponics-related example could be:

\begin{quote}
``Create a compact hydroponic support product that improves nutrient monitoring for small-scale growers and can be prototyped within one academic semester.''
\end{quote}

This second example is broad enough to allow several concepts, but narrow enough to keep the team aligned.

\section{Pre-Class Activity: Draft Candidate Charters}

Before class, each team member should draft at least two innovation charters. Each charter should be one or two sentences long.

\subsection{Drafting Guidance}

When writing your charter, try to include:

\begin{itemize}
    \item The general problem or opportunity area;
    \item The product category or type of offering;
    \item The target user or market context; and
    \item An appropriate time frame or development ambition.
\end{itemize}

Avoid detailed solution descriptions at this stage. If a charter already specifies the exact mechanism, geometry, or final concept, it is probably too narrow.

\subsection{Questions to Check Each Draft}
\begin{enumerate}
    \item Does the charter stay within the hydroponics theme?
    \item Does it allow multiple possible product opportunities?
    \item Does it leave room for applied physics analysis or technical justification?
    \item Is the scope ambitious but still manageable within the course?
\end{enumerate}

Bring your two best drafts to class with a one-line note explaining why each one is promising.

\section{In-Class Activity: Compare, Select, and Refine}

In class, your team should compare the individual charters and agree on one final version.

\subsection{Suggested Discussion Process}

%%%%%%%%%%%%%%%%%%%%%%%%ssss 





During this stage, it is expected that each group member has individually generated several innovation charters according to the given guidelines.
\begin{enumerate}
    \item Read all candidate charters aloud;
    \item identify the strengths and limitations of each draft;
    \item compare them using the criteria in Table~\ref{tab:session3-charter-criteria};
    Discuss with all members whether it is more advantageous to select the best innovation charter that had been drafted by one of the members, or to combine elements from multiple charters into a concise sentence.
    \begin{enumerate}
        \item Consider the diverse perspectives of each group member during the discussion.
        \item Evaluate the benefits and drawbacks of choosing a single charter versus combining elements from multiple charters by take into account the shared vision, goals, and desired outcomes of the project.
    \end{enumerate}
    \item Once consensus is reached, draft a new innovation charter that reflects the collective agreement of all team members.
    \begin{enumerate}
        \item Ensure that the new innovation charter is clear, concise, and representative of the team vision and objectives.
        \item Review and refine the drafted innovation charter to ensure that it aligns with the desired outcomes and sets appropriate boundaries for the innovation effort.
    \end{enumerate}
\end{enumerate}
%%%%%


\begin{table}[htbp]
\centering
\caption{Criteria for reviewing candidate innovation charters}
\label{tab:session3-charter-criteria}
\renewcommand{\arraystretch}{1.2}
\begin{tabularx}{\textwidth}{|>{\raggedright\arraybackslash}p{0.2\textwidth}|X|}
\hline
\textbf{Criterion} & \textbf{Question to ask} \\
\hline
Relevance & Does the charter clearly fit hydroponic agriculture? \\
\hline
Breadth & Does the charter allow several opportunities rather than one fixed solution? \\
\hline
Feasibility & Can the resulting project be explored within one semester? \\
\hline
Physics connection & Will the charter support later analysis involving sensing, flow, structure, energy, or control? \\
\hline
Potential impact & Is the opportunity area meaningful for users or operators? \\
\hline
\end{tabularx}
\end{table}

\subsection{Expected Output}

At the end of the session, your group should have one final innovation charter and a short explanation of why it was selected.

\section{Next Steps}

Save the final charter in your group records and upload it if the lecturer requests a submission. In Session~\ref{ch:opportunity-identification}, you will use the charter as the basis for generating and screening specific innovation opportunities.
